% requer: longtable, booktabs, array
\subsubsection{Planeamento da viagem}

\begin{longtable}{@{}p{4cm}p{11.3cm}@{}}
\caption{Glossário do modelo de domínio.}\label{tab:glossario-planeamento}\\
\toprule \textbf{Termo} & \textbf{Significado} \\ \midrule \endfirsthead
\toprule \textbf{Termo} & \textbf{Significado} \\ \midrule \endhead
Viagem & A viagem que o grupo quer fazer: destino, datas possíveis e tudo o que se planeia à volta. É a raiz do contexto.  \\
Ideia & Uma ligação de vídeo ou um texto que dá início a uma viagem, guardada com a origem à vista. O sistema nunca interpreta o vídeo; pede o destino em texto.  \\
Orçamento & O limite por pessoa definido pelo organizador. Todos os preços mostrados são comparados com ele.  \\
Proposta & O que o grupo vê para decidir: as opções, os preços e o total por pessoa, partilhados numa só ligação. Reflete sempre as alterações posteriores.  \\
Opção & Uma coisa que pode entrar numa proposta (voo, alojamento ou atividade), com ligação de origem, preço e estado: candidata, cortada (riscada mas visível) ou acordada.  \\
Opção de voo / alojamento / atividade & Tipos de Opção. Um voo sabe se não há transporte para o aeroporto à hora de partida; um alojamento tem noites e preço por noite; uma atividade (do roteiro da IA) fica confirmada ou descartada depois de verificada no mapa.  \\
Preço & Preço base mais extras (p. ex. bagagem de porão), dando o preço final; se o site esconder os extras, aparecem como \guillemotleft{}por confirmar\guillemotright{} em vez de serem omitidos.  \\
Total por pessoa & Soma do preço final por pessoa de todas as opções não cortadas, mais o transporte até ao aeroporto. Avisa quando ultrapassa o Orçamento.  \\
\bottomrule
\end{longtable}

\subsubsection{Decisão de grupo}

\begin{longtable}{@{}p{4cm}p{11.3cm}@{}}
\caption{Glossário do domínio \textit{Decisão de grupo}}\label{tab:glossario-decisão}\\
\toprule \textbf{Termo} & \textbf{Significado} \\ \midrule \endfirsthead
\toprule \textbf{Termo} & \textbf{Significado} \\ \midrule \endhead
Grupo & As pessoas que viajam juntas. Tem exatamente um organizador.  \\
Membro & Uma pessoa do grupo que pode dar opiniões e indicar disponibilidade.  \\
Organizador & O membro que planeia: partilha a proposta, define o prazo, regista decisões e pagamentos. Só o organizador pode marcar um pagamento como feito.  \\
Opinião (\guillemotleft{}Concordo\guillemotright{} / \guillemotleft{}Prefiro outra\guillemotright{}) & A posição de um membro sobre uma opção. O organizador vê em que itens há divergência antes de sair do Roamly para reservar.  \\
Disponibilidade & Se um membro pode nas datas propostas, ou em que outras datas pode.  \\
Decisão & O que o grupo combinou, com um prazo enquanto está aberta. Fica junto das opções a que se refere.  \\
\bottomrule
\end{longtable}

\subsubsection{Vigilância de preços}

\begin{longtable}{@{}p{4cm}p{11.3cm}@{}}
\caption{Glossário do domínio \textit{Vigilância de preços}}\label{tab:glossario-vigilância}\\
\toprule \textbf{Termo} & \textbf{Significado} \\ \midrule \endfirsthead
\toprule \textbf{Termo} & \textbf{Significado} \\ \midrule \endhead
Oferta vigiada & Uma ligação guardada que o Roamly volta a ler continuamente. Se não a conseguir ler, o organizador preenche o preço à mão.  \\
Leitura de preço & Um preço lido num dado momento.  \\
Alerta de preço & Gerado quando o preço de uma oferta vigiada sobe: mostra de quanto e notifica todos os membros do grupo, haja ou não um prazo pendente. O Roamly não congela preços nem reserva.  \\
\bottomrule
\end{longtable}

\subsubsection{Acerto de pagamentos}

\begin{longtable}{@{}p{4cm}p{11.3cm}@{}}
\caption{Glossário do domínio \textit{Acerto de pagamentos}}\label{tab:glossario-pagamentos}\\
\toprule \textbf{Termo} & \textbf{Significado} \\ \midrule \endfirsthead
\toprule \textbf{Termo} & \textbf{Significado} \\ \midrule \endhead
Acerto & Tudo o que é devido numa viagem depois de as opções serem acordadas, incluindo o que o organizador tem de adiantar.  \\
Parte da despesa & O que cada membro deve, calculado a partir das opções acordadas.  \\
Pagamento & Um registo de que um membro pagou a sua parte, feito pelo organizador. O Roamly regista; não processa dinheiro.  \\
Lembrete de pagamento & Enviado aos membros cuja parte não está marcada como paga, na altura que o organizador escolher.  \\
\bottomrule
\end{longtable}

\subsubsection{Balcão de agência}

\begin{longtable}{@{}p{4cm}p{11.3cm}@{}}
\caption{Glossário do domínio \textit{Balcão de agência}}\label{tab:glossario-agência}\\
\toprule \textbf{Termo} & \textbf{Significado} \\ \midrule \endfirsthead
\toprule \textbf{Termo} & \textbf{Significado} \\ \midrule \endhead
Quadro do cliente & Um quadro por cliente, com os pedidos pendentes e as propostas por enviar.  \\
Pedido do cliente & O roteiro que o cliente traz, o orçamento, as datas flexíveis e \guillemotleft{}o que não abdicam\guillemotright{}; assinala os campos em falta.  \\
Lembrete no quadro & Uma nota no quadro de um cliente (substitui o papel autocolante).  \\
Fornecedor & Fornecedor que não está online, guardado com as notas do agente; aparece entre as opções.  \\
Revisão da proposta & A assinatura do agente; o cliente não vê nada antes dela.  \\
\bottomrule
\end{longtable}